\section{系统实现与关键文件说明}

为方便说明项目整体结构，本项目将主要代码分为文献抓取与信息抽取、图数据库导入、前端可视化展示以及智能问答后端四个部分。各关键文件作用如下表所示。

\begin{longtable}{p{4cm} p{9cm}}
  \toprule
  文件/模块 & 说明 \\
  \midrule
  index\_demo.html & 图谱展示与问答前端页面 \\
  api/qa.php & 智能问答后端接口 \\
  llm\_relay.py & 本地大模型中转服务 \\
  normalize\_line1\_graph.py & 数据清洗与命名规范化脚本 \\
  import\_*.cypher & Neo4j 节点与关系导入脚本 \\
  output.jsonl / normalized\_output.jsonl & 文献信息抽取结果与规范化结果 \\
  \bottomrule
\end{longtable}

\section{小组成员分工}

\begin{table}[H]
  \centering
  \caption{小组成员分工}
  \renewcommand{\arraystretch}{1.35}
  \begin{tabular}{p{2.3cm} p{10.5cm}}
    \toprule
    成员 & 主要贡献 \\
    \midrule
    \multirow{4}{*}{甲}
      & 负责 PubMed 文献检索与 LINE-1 相关摘要数据整理。 \\
      & 负责摘要级实体与关系的初步抽取，以及原始结构化结果输出。 \\
      & 负责提供文本挖掘模块的关键词设计、抽取样例与评估材料。 \\
      & 参与文献挖掘结果核对，并协助完善实验报告相关内容。 \\
    \midrule
    \multirow{4}{*}{乙}
      & 负责 Neo4j 图数据库建模、节点关系导入与数据规范化处理。 \\
      & 负责知识图谱可视化前端页面开发、双语展示与交互效果调整。 \\
      & 负责基于本地知识库和大模型接口的智能问答模块实现与联调。 \\
      & 负责系统集成测试、演示流程整理及实验报告主体撰写。 \\
    \bottomrule
  \end{tabular}
\end{table}

\section{总结与不足}

目前，本项目已经实现了从文献获取、信息抽取、图数据库建模、图谱可视化展示到基于本地知识库的智能问答的完整闭环。系统能够围绕 LINE-1 及其亚家族构建知识子图，支持疾病、功能和文献证据的图谱化展示，并可通过自然语言问题从本地知识图谱中检索证据，再结合大模型生成带参考文献的回答。

项目当前仍存在若干局限。首先，图谱中的部分实体名称仍存在中英文混杂和粒度不完全统一的问题，后续仍可继续规范化。其次，当前前端以演示型子图为主，若扩展到更大范围的 TE 类型，还需要进一步优化动态加载与布局策略。最后，问答模块虽然已经支持模板化查询和受限的 Text2Cypher 思路，但在更复杂、开放的问题场景下仍有提升空间。

总体而言，本项目已经达到课程大作业的核心目标，并为后续扩展到更大规模的 TE 文献知识图谱提供了可复用的技术框架。

\section*{正式写作前的材料准备清单}
\begin{itemize}[leftmargin=2em]
  \item 文献抓取与信息抽取脚本名称、运行流程与关键词策略
  \item 实体类型与关系类型清单
  \item Neo4j 节点与关系设计表
  \item RAG 架构图
  \item Prompt 设计说明
  \item 3 到 5 张系统截图
  \item 2 到 3 条典型问答示例
  \item 人工抽样准确率评估结果
  \item 小组成员分工说明
  \item GitHub 仓库地址与主要文件说明
\end{itemize}
